\section{A width-independent positive cold-scale range}
\label{app:spectral-cold}

An operator-norm truncation separates a small Gaussian background from
a bounded-energy network. The background need not have small Frobenius
norm: its operator norm controls its effect at every bounded virtual
hidden state. The excess above that operator-norm threshold has Gaussian
concentration independent of learner width. This gives a positive range
of scales on which the stronger energy-based sample bounds hold at
\emph{all} learner widths, rather than requiring $t^2\le1/N$.

\subsection{Gaussian spectral excess and its posterior tail}

For a square matrix $W$ and $a\ge0$, let $\mathcal B_a$ be the closed
operator-norm ball of radius $a$. If the singular values of $W$ are
$\sigma_i$, its Frobenius projection onto $\mathcal B_a$ clamps each
singular value to $\min\{\sigma_i,a\}$. The residual therefore has
singular values $(\sigma_i-a)_+$. Fix the absolute threshold constant
\[
 a_0=8.
\]
For a network $w$, leave the first layer and readout unchanged and
replace each internal matrix by its residual above $a_0t$:
\[
 V^1=W^1,\qquad v_V=v,\qquad
 V^s=W^s-B^s,\quad
 B^s=\operatorname{Proj}_{\mathcal B_{a_0t}}(W^s)
              \quad(2\le s\le H).
\]
In particular $\|B^s\|_{\rm op}\le a_0t$.
Define the spectral-excess radius
\begin{equation}\label{eq:spectral-radius}
 \mathcal E_t(w)^2=\kappa_N(V)
 =\frac dN\|W^1\|_F^2+N\|v\|_2^2
       +\sum_{s=2}^H
            \operatorname{dist}_F(W^s,\mathcal B_{a_0t})^2.
\end{equation}

\begin{lemma}[A width-independent spectral-excess estimate]
\label{lem:spectral-excess}
Under the raw Gaussian law $\mu_{N,t}$, with $u=t^2$,
\[
 \mathbb E\mathcal E_t^2\le u(d+2H),
\]
and for every $s\ge0$,
\begin{equation}\label{eq:spectral-concentration}
 \mu_{N,t}\{\mathcal E_t\ge\mathbb E\mathcal E_t+s\}
       \le\exp\left(-\frac{Ns^2}{2u}\right).
\end{equation}
\end{lemma}
\begin{proof}
A pair of $1/4$-nets of the unit sphere, each of cardinality at most
$9^N$, gives for an internal Gaussian matrix the elementary bound
\[
 \Pr\{\|W\|_{\rm op}>x\}
 \le2\exp\left(2N\log9-\frac{Nx^2}{8u}\right).
\]
Indeed, the operator norm is at most twice the largest bilinear form
on the two nets, and each such form is Gaussian of variance $u/N$.
Writing $R=\operatorname{dist}_F(W,\mathcal B_{8t})$, we have
$R^2\le N(\|W\|_{\rm op}-8t)_+^2$. Hence
\begin{align*}
 \mathbb E R^2
 &\le 4N\int_0^\infty s
       \exp\left(2N\log9-
                   \frac{N(8t+s)^2}{8u}\right)ds\\
 &\le 16u\exp[-N(8-2\log9)]\le u.
\end{align*}
The last inequality follows from $N\ge1$ and
$16\exp[-(8-2\log9)]<1$.
The expected first-layer contribution to \eqref{eq:spectral-radius}
is $u(d+1)$ and the readout contribution is $u$.
There are $H-1$ internal matrices. Thus the expectation is at most
$u(d+H+1)\le u(d+2H)$.

In the weighted coordinates
\[
 y=\bigl(\sqrt{d/N}\,W^1,\ W^2,\ldots,W^H,\ \sqrt N\,v\bigr),
\]
the raw law has independent coordinates of common variance $u/N$.
The quantity $\mathcal E_t$ is Euclidean distance to the closed convex
product set whose first and readout coordinates are zero and whose
internal blocks belong to $\mathcal B_{8t}$. Distance to a nonempty
closed set is 1-Lipschitz. Thus $\mathcal E_t$, as a function of a
standard Gaussian vector, is $t/\sqrt N$-Lipschitz.

For completeness, the Gaussian Lipschitz concentration estimate used
here follows from the log-Sobolev inequality and Herbst's calculation.
For a positive smooth function $g$, Ornstein--Uhlenbeck entropy
dissipation and score commutation give
\[
 \operatorname{Ent}_{\gamma}(g)
 \le\frac12\mathbb E_\gamma\frac{\|\nabla g\|^2}{g}.
\]
For $g=e^{\lambda F}$ and a function $F$ of Lipschitz constant $\sigma$,
this implies
$\log\mathbb E e^{\lambda(F-\mathbb EF)}\le\lambda^2\sigma^2/2$:
if $\chi(\lambda)=\log\mathbb Ee^{\lambda F}$, the entropy inequality
says $\lambda\chi'-\chi\le\lambda^2\sigma^2/2$, and integration
from zero proves the assertion. Markov's inequality and optimization in
$\lambda$ give the upper tail $e^{-s^2/(2\sigma^2)}$.
Smooth approximation gives the same result for Lipschitz distance
functions. Substituting $\sigma=t/\sqrt N$ proves
\eqref{eq:spectral-concentration}.
\end{proof}

Let $D_*=D_*^{\rm prod}$ be the product target-cylinder cost. The actual
posterior, for every correctly labeled training sample, obeys
\[
 Q_{S,N,t,c}\le\exp(ND_*/u)\mu_{N,t}.
\]
Fix $0<\varepsilon<1$ and set
\begin{equation}\label{eq:spectral-budget}
 \begin{split}
 \beta&=\varepsilon/2,\\
 K_{\rm sp}&=2(d+2H)+4[D_*+\log(1/\beta)],\\
 r_{\rm sp}&=\min\{1/6,c/\sqrt{1+HK_{\rm sp}}\},
 \qquad \rho_{\rm sp}=r_{\rm sp}/2.
 \end{split}
\end{equation}
Then, uniformly over all $N\ge C$ and $0<u\le1$,
\begin{equation}\label{eq:spectral-posterior-tail}
 Q_{S,N,t,c}\{\kappa_N(V)>K_{\rm sp}\}\le\beta.
\end{equation}
To see this, put $b_0=d+2H$ and $A=D_*+\log(1/\beta)$.
The inequality
$\sqrt{2b_0+4A}\ge\sqrt{b_0}+\sqrt{2A}$ and the preceding lemma imply
\[
 \mu_{N,t}\{\mathcal E_t^2>K_{\rm sp}\}
 \le \exp(-NA/u).
\]
Multiply by the posterior density bound. Since $N/u\ge1$, the result
is at most $\exp[-N\log(1/\beta)/u]\le\beta$.

\subsection{Combining spectral truncation with an energy code}

On the retained event, the original readout satisfies
$E_v=N\|v\|^2\le\kappa_N(V)\le K_{\rm sp}$.
Every original network in the posterior support is therefore robust on
the common radius $r_{\rm sp}$. No robustness of the residual network
$V$ is required.

Apply either the dense sampling code or the multilevel code to $V$,
with energy cutoff $K_{\rm sp}$ and disturbance budget
$\rho_{\rm sp}=r_{\rm sp}/2$. Both constructions have full virtual
hidden states $\widehat h^s(x)\in[-1,1]^N$ and reproduce their decoded
output by injections $\eta$ into $V$, with
$\|\eta\|_*\le\rho_{\rm sp}$.
For the dense construction the proved bound is even
$\rho_{\rm sp}/2$; the weaker common bound suffices here.
At these \emph{same} virtual states define injections into the original
network by
\[
 \xi^1=\eta^1,\qquad
 \xi^s=\eta^s+
   \tanh(V^s\widehat h^{s-1})-
   \tanh(W^s\widehat h^{s-1})\quad(2\le s\le H),
 \qquad \xi^{\rm out}=\eta^{\rm out}.
\]
The readout is unchanged, and these injections reproduce exactly the
same states and decoded output. Since tanh is 1-Lipschitz,
\[
 \frac1{\sqrt N}\left\|
 \tanh(V^s\widehat h^{s-1})-
 \tanh(W^s\widehat h^{s-1})\right\|_2
 \le\|B^s\|_{\rm op}\le8t.
\]
Therefore
\[
 \|\xi\|_*\le\rho_{\rm sp}+8t\sqrt{H-1}.
\]
The following positive scale ceiling is sufficient:
\begin{equation}\label{eq:spectral-cutoff}
 u_c=\frac{r_{\rm sp}^2}{256H}.
\end{equation}
For $t^2\le u_c$, the right side is at most $r_{\rm sp}$.
The original support condition consequently gives clipped-score error
at most $1/16$ for the decoded representative at every Boolean input.
The representative and the residual network $V$ need not satisfy that
support condition themselves.

Let $J_{\rm den}(K,\rho)$ denote the dense-code budget in
\eqref{eq:budget}, recalculated with energy $K$ and radius $\rho$.
Let $J_{\rm ML}(K,\rho)$ denote the bound in
\eqref{eq:mlmc-budget} at those same resources. Put
\begin{equation}\label{eq:spectral-code-budget}
 J_{\rm sp}=\min\{J_{\rm den}(K_{\rm sp},\rho_{\rm sp}),
                    J_{\rm ML}(K_{\rm sp},\rho_{\rm sp})\}.
\end{equation}
Choose whichever finite comparison class has the smaller known upper
bound; this choice is made before the sample is drawn and needs no
extra union bound.

\begin{theorem}[Positive cold scales uniformly over learner width]
\label{thm:spectral-cold}
For the Boolean teacher, Gaussian tanh learner, and unchanged robustness
support of the paper, suppose
\[
 n\ge\frac2\varepsilon\,[J_{\rm sp}+\log(1/\delta)].
\]
Then, with sample probability at least $1-\delta$,
\[
 R(Q_{S,N,t,c})\le\varepsilon
\]
simultaneously for every $N\ge C$ and every $0<t^2\le u_c$, with
$K_{\rm sp},r_{\rm sp},u_c,J_{\rm sp}$ as above.
The ceiling $u_c$ is positive and independent of learner width.
The same conclusion holds with optional label-independent global-margin
conditioning.
\end{theorem}
\begin{proof}
For every retained original network that fits the training margins,
the preceding construction supplies a representative in the common
chosen finite class. Its clipped score differs by at most $1/16$ at
every input. Its empirical ramp loss is therefore zero, and its ramp
loss is one at every test-margin failure of the original network.
For any fixed representative with mean ramp loss greater than
$\varepsilon/2$, the probability of zero empirical ramp loss is at
most $\exp(-n\varepsilon/2)$. A union bound over the finite class shows
that every fitting retained original network has test-margin error
at most $\varepsilon/2$ on one sample event of probability at least
$1-\delta$. This event is independent of $N,t$.
Adding the discarded posterior mass in
\eqref{eq:spectral-posterior-tail}, at most $\beta=\varepsilon/2$,
proves the claim. The optional initial global-margin condition keeps
the same target-cylinder denominator and the same support, so all
steps remain valid.
\end{proof}

At fixed $d,c,\varepsilon,\delta$, the product-cylinder calculation gives
\[
 K_{\rm sp}=O(CLF_0),\qquad
 r_{\rm sp}^{-2}=O(CL^2F_0),\qquad
 u_c^{-1}=O(CL^3F_0),
 \qquad F_0=1+\log^2(2+C).
\]
Thus the same sufficient sample order as in the stronger-scale energy
result holds throughout this width-independent cold interval:
\[
 \widetilde O\!\left(\min\{C^4L^7,C^{9/2}L^6\}\right).
\]
Its log-retaining form, with $G=\log(2+CL)$, is
\[
 O\!\left(G\min\left\{
 C^4L^7F_0^4\log^2(2+L),\quad
 C^{9/2}L^6F_0^{9/2}\log(2+L)
 \right\}\right).
\]
The finite-scale statement is stronger than merely taking a limit at
each fixed width: it provides a positive scale threshold common to
all widths above target complexity.
